\subsection{Présentation du dataset}

Nous utilisons le dataset \textbf{Sentiment140}, disponible sur HuggingFace Datasets.
Il contient \textbf{1 600 000 tweets} collectés via l'API Twitter entre avril et juin 2009,
annotés automatiquement par la polarité des émoticônes présentes dans les messages.

\begin{table}[H]
\centering
\caption{Caractéristiques du dataset Sentiment140}
\begin{tabular}{ll}
\toprule
\textbf{Propriété} & \textbf{Valeur} \\
\midrule
Nombre de tweets   & 1 600 000 \\
Classes            & Négatif (0), Positif (1) \\
Distribution       & 800 000 / 800 000 (parfaitement équilibré) \\
Longueur moyenne   & $\approx$ 72 caractères \\
Nombre de mots moy.& $\approx$ 12 mots \\
Langue             & Anglais \\
Source             & API Twitter (2009) \\
Annotation         & Automatique (émoticônes) \\
\bottomrule
\end{tabular}
\end{table}

\subsection{Analyse exploratoire}

\begin{figure}[H]
\centering
\includegraphics[width=\textwidth]{figures/01_data_understanding.png}
\caption{Analyse exploratoire complète du corpus Sentiment140 : distribution des classes,
longueur des tweets, nombre de mots, éléments de bruit détectés et statistiques générales.}
\label{fig:eda}
\end{figure}

\subsection{Problèmes potentiels identifiés}

L'analyse exploratoire révèle plusieurs défis spécifiques aux données Twitter :

\begin{itemize}
  \item \textbf{Bruit textuel} : présence massive d'URLs, @mentions, \#hashtags et emojis.
  \item \textbf{Langue informelle} : abréviations (\textit{lol, omg, tbh}), contractions (\textit{I'm, don't}).
  \item \textbf{Longueur très courte} : contrainte des 140 caractères, peu de contexte disponible.
  \item \textbf{Classes équilibrées} : le dataset est parfaitement équilibré (50\%/50\%), ce qui est favorable.
  \item \textbf{Annotation automatique} : quelques erreurs d'annotation possibles via les émoticônes.
\end{itemize}
